\chapter{Contexto y Motivaci\'on}

El proyecto ``Astro Denoise'' nace de la necesidad de evaluar, con rigor cient\'ifico, c\'omo las t\'ecnicas modernas de restauraci\'on de im\'agenes afectan la capacidad de las encuestas astron\'omicas para descubrir y caracterizar objetos tenues. El Vera Rubin Observatory (LSST) representa una revoluci\'on en volumen y profundidad de datos: detectar señales d\'ebiles en im\'agenes coadded es cr\'itico para objetivos como cosmolog\'ia de precisi\'on, mapeo de la V\'ia L\'actea, descubrimiento de objetos transitorios y estudios del Sistema Solar. Mejorar la relaci\'on señal\textendash ruido sin introducir sesgos es, por tanto, una prioridad cient\'ifica.

Este proyecto integra dos dimensiones esenciales: (1) rigor metodol\'ogico \textemdash protocolos reproducibles, m\'etricas cient\'ificas y control de semillas/random states\textemdash y (2) relevancia cient\'ifica \textemdash medir el impacto de la reducci\'on de ruido sobre la capacidad de detecci\'on y la calidad fotom\'etrica en cat\'alogos que alimentan an\'alisis astrof\'isicos. En otras palabras, no se persigue la m\'axima mejora perceptual, sino la mejora cuantificable y verificable en m\'etricas que importan para la ciencia.

La elecci\'on del conjunto de datos (DC2 dr6\textendash v4) y la atenci\'on a reproducibilidad responden a las recomendaciones de la comunidad y a los objetivos cient\'ificos del Rubin Observatory \cite{objetivos_rubin_2024}, que enfatizan descubrimientos estad\'isticos, trazabilidad de an\'alisis y estimaci\'on rigurosa de incertidumbres instrumentales.

Motivaci\'on inspiracional: este trabajo pretende acelerar descubrimientos al reducir las barreras t\'ecnicas que impiden detectar y medir fuentes d\'ebiles. Al construir un benchmark modular, reproducible y abierto, se facilita la comparaci\'on objetiva entre m\'etodos cl\'asicos y de IA, promoviendo buenas pr\'acticas de ciencia de datos en astronom\'ia. El impacto potencial incluye: mejores cat\'alogos para cosmolog\'ia y astrof\'isica, pipelines reutilizables para futuras encuestas y un marco de evaluaci\'on que favorece transparencia, reproducibilidad y colaboraci\'on abierta.

Se resalta la importancia de documentar cada decisi\'on \textemdash desde criterios de selecci\'on bibliogr\'afica hasta preprocesado y par\'ametros del detector\textemdash para que los resultados sean auditables y replicables por cualquier equipo de la comunidad.
